\section{三种绘景}

Gross 第 14 章强调：可观测量是矩阵元，因而可用随时间变化的幺正变换在不同“绘景”之间移动时间依赖，而不改变任何测量结果。

\subsection{绘景变换的一般结构}

若 $A(t)$ 对每一时刻幺正，则
\begin{equation}
  |\psi\rangle_A=A(t)|\psi\rangle,
  \qquad
  O_A(t)=A(t)O A^\dagger(t)
\end{equation}
保持 $\langle\psi|O|\phi\rangle$ 不变。三种标准选择如下（取 $\hbar=1$，参考时刻 $t_0=0$）。

\subsection{薛定谔绘景}

$A=1$。态满足
\begin{equation}
  i\partial_t|\psi_S\rangle=H_S(t)|\psi_S\rangle,
\end{equation}
算符只在 $H$ 显含时间时才随 $t$ 变。时间演化算符
\begin{equation}
  U(t,t')
  =T\exp\Bigl(-i\int_{t'}^{t}\mathrm{d}\tau\,H_S(\tau)\Bigr)
\end{equation}
满足 $i\partial_t U=H_S(t)U$、$U(t',t')=1$。

\subsection{海森堡绘景}

取 $A(t)=U^\dagger(t,0)$。则态冻结为初态，$|\psi_H\rangle=|\psi_S(0)\rangle$，算符
\begin{equation}
  O_H(t)=U^\dagger(t,0)O_S U(t,0).
\end{equation}
对 $t$ 求导并用 $i\partial_t U=HU$ 得海森堡方程
\begin{equation}
  i\partial_t O_H
  =[O_H,H_H]
  +\bigl(\partial_t O_S\bigr)_H.
\end{equation}
格林函数通常在此绘景用海森堡场算符定义。

\subsection{相互作用（Dirac）绘景}

把 $H=H_0+V$（$H_0$ 不含时），令
\begin{equation}
  |\psi_I\rangle
  =e^{iH_0 t}|\psi_S\rangle,
  \qquad
  O_I(t)
  =e^{iH_0 t}O_S e^{-iH_0 t}.
\end{equation}
则 $H_0$ 在薛定谔与相互作用绘景中相同：$[H_0]_I=H_0$。态的方程为
\begin{equation}
  i\partial_t|\psi_I\rangle=V_I(t)|\psi_I\rangle,
  \qquad
  V_I(t)=e^{iH_0 t}V e^{-iH_0 t},
\end{equation}
算符则按自由 $H_0$ 运动：
\begin{equation}
  i\partial_t O_I=[O_I,H_0]+(\partial_t O_S)_I.
\end{equation}
相互作用绘景演化算符
\begin{equation}
  U_I(t,t')
  =T\exp\Bigl(-i\int_{t'}^{t}\mathrm{d}\tau\,V_I(\tau)\Bigr)
\label{eq:UI}
\end{equation}
是零温图展开与 Gell-Mann--Low 定理的出发点。

\subsection{产生湮灭算符在相互作用绘景中的相位}

设 $H_0=\sum_j\varepsilon_j c_j^\dagger c_j$。在任意占据数本征态 $|\phi_\alpha\rangle$ 上，
\begin{equation}
  H_0 c_j|\phi_\alpha\rangle
  =(\varepsilon_j\text{ 已从该态扣除后的能量})
  =(W_\alpha-\varepsilon_j)c_j|\phi_\alpha\rangle,
\end{equation}
因而
\begin{equation}
  c_j(t)_I
  =e^{iH_0 t}c_j e^{-iH_0 t}
  =c_j e^{-i\varepsilon_j t},
  \qquad
  c_j^\dagger(t)_I
  =c_j^\dagger e^{i\varepsilon_j t}.
\label{eq:c_I}
\end{equation}
（对 $|\phi_\alpha\rangle$ 逐态验证后由完备性推广到算符等式。）这是画 $G_0$ 线时“粒子线带 $e^{-i\xi t}$、空穴线带 $e^{-i|\xi|t}$”的来源。

\begin{note}
$H=H_0+V$ 的划分原则上任意，只要 $H_0$ 不含时。实践中 $H_0$ 取独立粒子（自由或 HF），$V$ 为剩余相互作用；否则 $V_I(t)$ 的 Wick 收缩不再对应简单的 $G_0$。
\end{note}
